\scanpage{254}

et

\[\left\{\begin{aligned}
B_{0G}\cap G_0&=\unsure{\mathfrak{z}|\mathfrak{Z}}.B_{0H}\\
B_{1G}\cap G_0&=\mathfrak{z}.B_{1H}
\end{aligned}\right.\qquad\text{produits semi-directs},\tag*{(\unsure{112|111})}\]

formules qui proviennent de la conjonction de (108), et du fait que par définition

\[G_0=\operatorname{Ker}(G\xrightarrow{\delta}Z)\]

et qu'enfin $\delta$ s'explicite comme le composé

\[G=T_i.H\xrightarrow{\mathrm{can}}T_i\xrightarrow{\delta|T_i}Z\qquad i\in\{0,1\}\]

et \unsure{qu'|qui|que} \deleted{\unsure{en}}enfin $\deleted{\unsure{T_i\cap T_H}}\mathfrak{z}$ est le noyau de $\delta|T_i$.

Une partie des relations explicitées dans les formules précédentes sont résumées dans le diagramme d'inclusions

\diagram{\resizebox{\linewidth}{!}{\begin{tikzcd}[ampersand replacement=\&,column sep=small,row sep=large]
\& U_0 \arrow[r,hook,no head,"Z"] \arrow[d,hook,swap,"\mathfrak{z}"] \& B_{0H}=T_H.U_0 \arrow[r,hook] \arrow[d,hook,swap,"\mathfrak{z}"] \& H \arrow[d,hook]\\
(113) \& \widetilde U_0=\mathfrak{z}.U_0 \arrow[r,hook,no head,"Z"] \arrow[d,hook,swap,"\widetilde Z"] \& \widetilde B_0=\mathfrak{z}.B_{0H}=(\mathfrak{z}\times T_H).U_0 \arrow[r,hook] \arrow[d,hook,"\widetilde Z"] \& G_0=\mathfrak{z}.H \arrow[d,hook,"\widetilde Z"]\\
\& B'_0=T_0.U_0 \arrow[r,hook,no head,"Z"] \& B_{0G}=T_0.B_{0H}=(T_0\times T_H).U_0 \arrow[r,hook] \& G=T_0.H
\end{tikzcd}}}

dont \deleted{\unsure{les}} quatre carrés composants sont cartésiens et cocartésiens, et où on a écrit, au dessus des inclusions distinguées, la valeur à isom près du quotient correspondant. J'ai posé

\[\widetilde Z\overset{\text{déf}}{=}G/H\underset{\mathrm{surj}.}{\overset{\unsure{c|e}}{\longrightarrow}}G/G_0\deleted{\unsure{\simeq}}\lhook\joinrel\xrightarrow{\delta} Z\tag*{(114)}\]

% margin note written sideways
\margin{on désigne par $\unsure{\widetilde\delta|\delta}:G\longrightarrow G/H\simeq\widetilde Z$ l'appl. can. de passage au quotient.}

d'où un hom can. composé \deleted{\unsure{de}} de noyau $\mathfrak{z}$

\[0\longrightarrow\mathfrak{z}\longrightarrow\widetilde Z\xrightarrow{c}Z\tag*{(115)}\]

% archive sheet number: 253
